% Compilen desde esta carpeta o desde la raíz del repositorio.
\makeatletter
\def\input@path{{../../preamble/}{book2026-2027/preamble/}}
\makeatother
\input{preamble.tex}

\usepackage[spanish]{babel}
\usepackage{lmodern}
\usepackage{titlesec}
\geometry{a4paper,margin=1.6cm}
\linespread{1.10}

\definecolor{HeadingColor}{HTML}{243B53}
\definecolor{InstructionColor}{HTML}{704214}
\definecolor{ConceptColor}{HTML}{205C50}
\definecolor{LinkColor}{HTML}{0047AB}
\hypersetup{allcolors=LinkColor}
\newcommand{\guidelink}[2]{\href{#1}{\textbf{\underline{#2}}}}
\newcommand{\instruction}[1]{\textcolor{InstructionColor}{\textbf{#1}}}
\newcommand{\concept}[1]{\textcolor{ConceptColor}{\textbf{#1}}}
\newcommand{\checkbox}{%
    \begingroup
    \setlength{\fboxrule}{0.4pt}%
    \setlength{\fboxsep}{2pt}%
    \fbox{\rule{0pt}{0.8ex}\rule{0.8ex}{0pt}}%
    \endgroup}

\setlength{\parskip}{4pt}
\setlength{\tabcolsep}{5pt}
\renewcommand{\arraystretch}{1.15}
\setlist{leftmargin=0.65cm,topsep=4pt,partopsep=0pt,parsep=0pt,itemsep=3pt}
\titleformat{\section}{\large\bfseries\color{HeadingColor}}{}{0pt}{}
\titlespacing*{\section}{0pt}{10pt}{3pt}
\newcolumntype{Y}{>{\raggedright\arraybackslash}X}

\hypersetup{pdftitle={C7--C8. Prueben el prototipo y ajusten la solución},pdfsubject={Emprendimiento e Innovación, Grado 10, Periodo 2}}

\begin{document}

\begin{center}
{\Large\bfseries\color{HeadingColor} C7--C8. Prueben el prototipo y ajusten la solución\par}
\vspace{3pt}
Emprendimiento e Innovación \enspace|\enspace Grado 10 \enspace|\enspace Periodo 2\\
Guía 4 de 5
\end{center}

\section*{Objetivo y punto de partida}
\concept{C7.} Prueba el prototipo con usuarios, pares o docentes para recoger impresiones iniciales.\\
\concept{C8.} Ajusta la solución a partir de retroalimentación recibida y de las limitaciones identificadas.

Recuperen V2, sus señales de uso y PL1 actualizado (P2, C5--C6). Consulten D1 (P1, C7) y el supuesto que buscan aclarar. Dediquen la prueba a una acción que puedan observar.

\section*{Documento acumulativo e integridad de la evidencia}
Esta guía es de consulta. Continúen el \textbf{mismo documento del proyecto en LaTeX} en \guidelink{https://www.overleaf.com/}{Overleaf} o en su editor habitual. Actualicen la fecha y entreguen el PDF completo. Usen \textbf{P1} para el Periodo 1 y \textbf{P2} para el Periodo 2 al citar criterios. Enlacen los registros por código o referencia interna y conserven sus fuentes en la sección de origen.

Distingan \concept{evidencia} (lo observado o expresado), \concept{conclusión} (lo que infieren), \concept{supuesto} (lo pendiente de verificar) e \concept{idea} (una propuesta). Registren fuentes y fechas reales. Identifiquen toda simulación autorizada por el docente. \instruction{Si usan IA, limítenla al código LaTeX, al formato y a organizar contenido propio.} Evidencias, conclusiones, ideas, decisiones y aportes deben provenir de su trabajo.

\section*{C7 de P2. Preparen, realicen y registren una prueba breve}
\begin{enumerate}[series=pasos]
\item \textbf{Definan la prueba PR1.} Identifiquen la versión, el contexto y a quién invitarán (usuarios, pares o docentes). Escriban una consigna que indique qué debe lograr la persona. Acuerden qué resultado o dificultad observarán, cuánto tiempo darán y quién facilitará y registrará.
\item \textbf{Realicen la prueba.} Presenten la misma consigna a los participantes sin indicarles los pasos. Si intervienen o simulan una parte, dejen constancia. Al terminar, pregunten qué les resultó fácil o difícil y en qué momento. Registren fecha y tipo de participante con un código, sin datos personales innecesarios.
\item \textbf{Separen las tres clases de registro.} Asignen un código a cada fila (por ejemplo, PR1-01). Describan la acción que vieron, anoten la respuesta como cita exacta o paráfrasis identificada y expliquen por separado qué interpretación propone el equipo.
\end{enumerate}

{\small
\begin{tabularx}{\linewidth}{|Y|Y|Y|}
\hline
\textbf{Lo que ocurrió} & \textbf{Lo que expresó la persona} & \textbf{Lo que interpretó el equipo} \\
\hline
Acción, resultado, dificultad o ayuda recibida. & Comentario registrado, con su contexto. & Posible explicación y duda pendiente. \\
\hline
\end{tabularx}\par
}

\concept{Evidencia de C7.} Protocolo breve y registros de la prueba realizada, con la tarea, la versión, el tipo de participante y las tres clases de registro diferenciadas.

\newpage
\section*{C8 de P2. Documenten qué cambió y por qué}
\begin{enumerate}[resume=pasos]
\item \textbf{Elijan ajustes que puedan realizar.} Revisen las filas de PR1 y las limitaciones de materiales, tiempo, acceso o alcance. Prioricen cambios según su efecto probable sobre la tarea y los recursos de PL1. Si dejan una sugerencia pendiente, indiquen por qué.
\item \textbf{Apliquen los cambios y conserven ambas versiones.} Identifiquen la versión ajustada como V3. Para cada ajuste, muestren el antes y el después con una imagen, captura o secuencia anotada. Citen la observación, el comentario o la limitación concreta que motivó el cambio.
\end{enumerate}

{\small
\begin{tabularx}{\linewidth}{|Y|Y|Y|Y|}
\hline
\textbf{Evidencia o límite} & \textbf{Antes (V2)} & \textbf{Después (V3)} & \textbf{Efecto esperado y comprobación} \\
\hline
Código PR1 o limitación concreta con su fuente. & Elemento o paso que generó la dificultad. & Cambio visible y ubicación en V3. & Qué debería mejorar y qué observarían para comprobarlo. \\
\hline
\end{tabularx}\par
}

\textbf{Ejemplo ficticio.} Si una persona se dirige a la fila de pago al intentar recoger un pedido, podrían mover la señal de recogida a un punto visible antes de que deba elegir una fila. Al repetir la tarea, observarían hacia dónde se dirige primero.

\begin{enumerate}[resume=pasos]
\item \textbf{Aclaren el alcance de lo aprendido.} Si repiten la tarea con V3, registren ese resultado por separado, con fecha y versión. Si queda por comprobar, señalen el efecto como esperado. Indiquen qué límites tuvo PR1 y actualicen PL1 y la bitácora con los acuerdos y pendientes que dejó la prueba.
\end{enumerate}

\concept{Evidencia de C8.} Una versión ajustada y un registro de cambios que conecta cada diferencia visible entre versiones con retroalimentación o limitaciones identificadas.

\section*{Resultado y continuidad}
Agreguen PR1, V3 y el registro de cambios. En P2, C9--C10 recuperarán esas referencias para explicar qué aprendieron juntos y qué respaldo tiene la mejora que presentarán en la feria.

\section*{Verificación final}
\begin{itemize}[label=\checkbox]
\item Realizamos una prueba con una tarea observable y registramos la versión usada.
\item Separamos lo ocurrido, lo expresado y nuestra interpretación.
\item Cada ajuste tiene un antes y un después y cita su motivo.
\item Actualizamos la bitácora y compilamos el PDF acumulativo con sus referencias.
\end{itemize}

\end{document}
